\begin{textsample}{POS Dim 8 – global\_north\_2025\_01 – Score 13.00 – global\_north\_2025\_01/120437501.txt}  \label{ex:f8_pos_106}
Palestinians returning to north Gaza, which has endured a devastating siege for months, are finding that their homes have been turned to rubble.
The area has been reeling following an Israeli military plan that legal experts have described to Middle East Eye as"genocidal"and aimed at the permanent forced \textbf{displacement} of northern Gaza ' s population, as settlers plan for post-war colonisation of the Palestinian enclave.
Proposed and promoted by a group of senior Israeli army reservists, the"Generals ' Plan", also known as the"Eiland Plan", involves forcibly displacing the entire population of northern Gaza, and then besieging the area, including stopping the entry of humanitarian supplies, in order to starve out anyone left.
The scheme has so far lead to the killing or disappearance of at least 5,000 people in the northern part of the enclave, with another 9,500 wounded as a result of the brutal campaign that began in early October last year, a medical source told Al Jazeera. \textbf{MEE} correspondents @ @ @ @ @ @ @ @ @ @ forcing unarmed and starving civilians in the north from their homes and \textbf{displacement} shelters at gunpoint.
New \textbf{MEE} \textbf{newsletter} : Jerusalem Dispatch Sign up to get the latest \textbf{insights} and analysis on Israel-Palestine, alongside Turkey \textbf{Unpacked} and other \textbf{MEE} \textbf{newsletters} Fadi, a resident of north Gaza, told \textbf{MEE} that he and his family returned after their expulsion by the Israeli army to find their home completely destroyed.
Asked if he had somewhere else to stay, he pointed at a tent he had taken with him from southern Gaza, saying it would take too long to wait for aid and shelter to be provided.
Fadi and 14 others from his family are expected to live in the small tent."My message to the world is that we want nothing from you,"Fadi said."We have spent around one year and three months resilient in the Gaza Strip, we did not need anything from the world, we only had God with us,"he added, recalling the lack @ @ @ @ @ @ @ @ @ @ Doctors Without Borders ( MSF ), the Israeli offensive in the north has left thousands of people without"access to food, water, or healthcare".
There have been no hospitals functioning in the area since 9 January.
As the north deals with a severe lack of basic necessities, Muhammad Dahr, a resident in northern Gaza who has been displaced for 14 months, said that he has hope that sufficient aid will enter into the enclave."Today we can't live, but God is standing by us and we pray for the best."Grief and loss in Gaza When speaking about whether the war will continue, Fadi said that is was a matter for"God ' s will", but that he would remain in the north despite his and his community ' s suffering."Our pain is deep,"he said, lamenting the loss of friends, loved ones and homes.
Dahr felt similarly as he spoke of his home and the loved @ @ @ @ @ @ @ @ @ @ war."My home was gone on the seventh day of the war... but when I saw all of Gaza ' s homes gone, I surrendered to God,"Dahr told \textbf{MEE}."Nothing remains...
All of Gaza is in ruins, even the houses of worship, no stone, no trees, nothing,"he said, adding that he did not even recognise his home or neighbourhood, such was the level of destruction.
Dahr told \textbf{MEE} that he had lost an \textbf{extraordinary} 600 members of his family.
Mosab Muhammad Saleh, a child who lost his father, told \textbf{MEE} that he had hoped they would still have their home intact.
But he and his remaining family found it completely demolished, and they now have to shelter in a nearby school.
When asked about his message to the world, the boy simply began to sing about his love of Palestine.
More than 47,000 Palestinians have been killed by Israeli forces in all of Gaza over the @ @ @ @ @ @ @ @ @ @.
Human rights groups and UN experts have accused Israel of collective punishment against Palestinians since the Hamas-led attacks of 7 October 2023.
One elderly woman, from the Saftawy family, recalled her \textbf{displacement} from northern Gaza as Israeli shells rained down.
When the ceasefire came into effect, she headed back towards her home.
But when they came back, their once four-storey family house was nothing more than a pile of rubble."I was expecting to at least find the lower floors of the house so that my daughters could shelter there, however I found it all in ruins, just rubble, as if it was crumbled up biscuits,"she said, adding that the damage was so great, that they could not even retrieve personal items such as clothing.
And yet, in the devastation, she - like many others - reached for resilience."Praise to God, we remain standing."Middle East Eye delivers independent and \textbf{unrivalled} \textbf{coverage} and analysis of the Middle @ @ @ @ @ @ @ @ @ @ about \textbf{republishing} this content and the \textbf{associated} \textbf{fees}, \textbf{please} \textbf{fill} out this form.
More about \textbf{MEE} can be found here.

% matched lemmas: associate, coverage, displacement, extraordinary, fee, fill, insight, mee, newsletter, please, republish, unpack, unrivalled
\end{textsample}
